\subsectionaddtolist{ پروتکل مبتنی بر رمزنگاری نامتقارن}

\subsubsection*{آ}
این پروتکل برای {احراز اصالت کلاینت} طراحی شده تا کاربر A ثابت کند همان کسی است که ادعا می‌کند.
\begin{itemize}
    \item {پیام اول ($A\rightarrow B:[ID]_{A}$):} کاربر A شناسه خودش را برای B می‌فرستد تا اعلام حضور کند و بگوید می‌خواهد ارتباط را شروع کند.
    \item {پیام دوم ($B\rightarrow A:N_{B}$):} کاربر B یک چالش (Challenge) به صورت یک نانس تصادفی ($N_B$) می‌فرستد. این کار برای جلوگیری از حمله Replay است تا مطمئن شود پاسخ کلاینت زنده و در لحظه است.
    \item {پیام سوم ($A\rightarrow B:E(PrvKey_{A},N_{B})$):} کاربر A نانس دریافتی را با کلید خصوصی خودش رمز (در واقع امضا) می‌کند و پس می‌فرستد. 
\end{itemize}
 چون فقط A به کلید خصوصی‌اش ($PrvKey_A$) دسترسی دارد و نانس هم جدید بوده، B با رمزگشایی این پیام توسط کلید عمومی A و تطبیق آن با $N_B$ مطمئن می‌شود که در همین لحظه دارد با خود واقعی A صحبت می‌کند.

\medskip
\subsubsection*{ب}
این پروتکل در برابر {حمله جعل هویت از طریق سرور واسط} آسیب‌پذیر است. مشکل اینجاست که امضای A روی نانس، هیچ وابستگی به هویت مقصد (B) ندارد. 
اگر A بخواهد همزمان با یک سرور مخرب یا کاربر دیگر به نام C ارتباط برقرار کند، C می‌تواند نانسی که B برایش فرستاده را عینا به عنوان چالش برای A بفرستد. A آن را امضا می‌کند و به C می‌دهد. حالا C امضای A روی نانس B را دارد؛ آن را به B می‌فرستد و خودش را به جای A جا می‌زند، بدون اینکه کلید خصوصی A را داشته باشد.



\subsectionaddtolist{ پروتکل ایجاد کلید جلسه}

\subsubsection*{آ}
\begin{itemize}
    \item A فکر می‌کند کلید $K_{AB}^{\prime}$ را با B به اشتراک گذاشته، چون این کلید داخل یک پیام رمز شده با کلید بلندمدت $K_{AB}$ آمده است؛ کلیدی که از نظر A فقط خودش و B آن را بلد هستند.
    \item B فکر می‌کند کلید $K_{AB}^{\prime}$ را با A به اشتراک گذاشته، چون در پیام سوم، A توانسته نانسی را که B تعیین کرده بود ($N_A$ که در پیام دوم داخل متن رمز بود) استخراج کند و حالا آن نانس را با خود کلید جلسه جدید ($K_{AB}^{\prime}$) رمز کرده و پس فرستاده است. این یعنی A پیام دوم را باز کرده و کلید جدید را قبول دارد.
    \item A فکر می‌کند کلید $K_{AB}^{\prime}$ تازه است، چون خودش در پیام اول نانس $N_A$ را فرستاده بود و در پیام دوم، همان نانس را به صورت رمز شده پس گرفته است. این یعنی پیام دوم قدیمی نیست و در پاسخ به درخواست فعلی او تولید شده است.
    \item B فکر می‌کند کلید $K_{AB}^{\prime}$ تازه است، چون خودش نانس $N_A$ را از پیام اول گرفته، کلید جدید را کنارش گذاشته و فرستاده است. نانس $N_A$ برای B نقش شناسه نشست فعلی را دارد. (البته این فرض ضعف اصلی پروتکل است، چون نانس را A تولید کرده نه B ).
\end{itemize}

\medskip
\subsubsection*{ب}

چون پروتکل متقارن است و امکان اجرای معکوس (شروع از سمت B ) وجود دارد، مهاجم C بدون داشتن کلید بلندمدت $K_{AB}$ می‌تواند از خود A به عنوان یک Oracle استفاده کند تا پیام رمز شده‌ی مورد نظر را برایش تولید کند. 

 کاربر A قصد دارد یک نشست عادی با B شروع کند و نانس $N_A$ را می‌فرستد. مهاجم این پیام را Intercept می‌کند:
    $$ A \rightarrow C(B): A, N_A$$
    
 مهاجم پیام بالا را نگه می‌دارد و یک اجرای معکوس (پروتکل رفرش از سمت B ) با خود A باز می‌کند و همان نانس $N_A$ را به عنوان چالش به A می‌فرستد تا خودش را جای B جا بزند:
    $$C(B) \rightarrow A: B, N_A$$
    
 حالا A در نقش پاسخ‌دهنده به اجرای معکوس، یک کلید تازه مانند $K^*$ تولید کرده و آن را همراه با نانس $N_A$ با کلید $K_{AB}$ رمز می‌کند و می‌فرستد:
    $$ A \rightarrow C(B): E(K_{AB}, [N_A, K^*])$$
    
 مهاجم که به این رمز دست پیدا کرده، همین پیام را عینا به عنوان پیام دوم اجرای اول به A برمی‌گرداند:
    $$ C(B) \rightarrow A: E(K_{AB}, [N_A, K^*])$$
    
 از نظر A این پیام معتبر است، چون با کلید $K_{AB}$ رمز شده و حاوی همان نانس اولیه خودش ($N_A$) است. پس فکر می‌کند پیام از طرف B آمده و تاییدیه نهایی را با کلید جدید $K^*$ می‌فرستد:
    $$ A \rightarrow C(B): E(K^*, N_A)$$
    
 مهاجم این تاییدیه را هم عینا به اجرای معکوس برمی‌گرداند تا آن نشست را هم تکمیل کند:
    $$ C(B) \rightarrow A: E(K^*, N_A)$$

 در پایان این فرآیند، B اصلا در جریان ارتباط نبوده و روحش هم از این ماجرا خبر ندارد، اما A کاملا فریب خورده و فکر می‌کند کلید تازه $K^*$ را با B به اشتراک گذاشته است، در حالی که این کلید فقط بین A و مهاجم (که نقش آینه را بازی کرده) قرار دارد. دلیل اصلی موفقیت این حمله، مشخص نبودن نقش فرستنده و گیرنده (هویت‌ها) در داخل بدنه پیام رمز شده است.

\subsubsection*{ج}
بهترین کار این است که B هم یک نانس اختصاصی ($N_B$) تولید کند و در رمزگذاری پیام دوم مشارکت دهد، یا اینکه هویت‌ها صراحتا داخل پیام رمز شده بیایند:
$$ B\rightarrow A: E(K_{AB}, [N_{A}, N_B, K_{AB}^{\prime}])$$
$$ A\rightarrow B: E(K_{AB}^{\prime}, N_{B})$$
اینطوری B هم مطمئن می‌شود پیام سوم در پاسخ به نانس جدید خودش تولید شده است.



\subsectionaddtolist{ پروتکل پرداخت PayDroid}

\subsubsection*{آ}
خیر، این پروتکل {اصلا امن نیست}. مشکل بزرگ این است که پیام حاوی هویت گیرنده ($B$) و مبلغ ($X$) به صورت آشکار فرستاده می‌شود و امضای دیجیتال A فقط روی $X$ و $n$ قرار دارد: $E(PrvKey_{A},[X,n])$.

مهاجم (مثلا خود Bob یا یک شنودکننده مثل C ) پیام آلیس را در شبکه Intercept می‌کند. مهاجم هویت $B$ را در بخش آشکار پیام حذف کرده و شناسه خودش (مثلا $C$) را قرار می‌دهد:
$$A\rightarrow S: A, C, X, n, Cert_{A}, E(PrvKey_{A},[X,n])$$
وقتی پیام به سرور S می‌رسد، سرور امضا را با کلید عمومی آلیس چک می‌کند. امضا درست است (چون مبلغ X و نانس n تغییر نکرده‌اند). سرور هم بدون اینکه خبر داشته باشد، پول را به حساب C واریز می‌کند!

\medskip
\subsubsection*{ب}
برای حل این مشکل، باید هویت گیرنده ($B$) هم حتما داخل پکیج امضا قرار بگیرد تا کسی نتواند مقصد پول را در مسیر عوض کند:
$$A\rightarrow S: A, B, X, n, Cert_{A}, E(PrvKey_{A}, [B, X, n])$$



\subsectionaddtolist{ توزیع کلید با سرور مرکزی}

\subsubsection*{آ}
\begin{itemize}
    \item {خط ۱ ($A \rightarrow S: A, B$):} کاربر A به سرور اعلام می‌کند که من (A) می‌خواهم با B یک ارتباط امن داشته باشم، لطفا یک کلید مشترک برای ما بساز.
    \item {خط ۲ ($S \rightarrow A: E(K_A, R), E(K_B, R)$):} سرور کلید تصادفی $R$ را تولید می‌کند. یک نسخه از آن را با کلید مخفی A رمز می‌کند تا فقط A بفهمد، و نسخه دیگر را با کلید مخفی B رمز می‌کند و هر دو را به A می‌دهد.
    \item {خط ۳ ($A \rightarrow B: E(K_B, R), E(R, M)$):} کاربر A بخش مربوط به B را برایش می‌فرستد تا B هم به کلید $R$ برسد. همزمان پیغام $M$ را با کلید جلسه $R$ رمز کرده و برای B ارسال می‌کند.
\end{itemize}

\medskip
\subsubsection*{ب}



هکر قادر است به پیغام M دسترسی پیدا کند. از آنجایی که این پروتکل هیچ مکانیزمی برای احراز اصالت سرور (مانند امضای دیجیتال سرور) و هیچ نانسی برای تضمین تازگی ندارد، در برابر جعل هویت سرور و \lr{Replay Attack} آسیب‌پذیر است.

اگر هکر (C) از قبل یک کلید نشست قدیمی ($R_{old}$) و پیام رمز شده‌ی آن ($E(R_{old}, M_{old})$) را به همراه تیکت B (یعنی $E(K_B, R_{old})$) شنود و ذخیره کرده باشد، می‌تواند مراحل زیر را اجرا کند:


 کاربر A طبق روال درخواست خود را به سمت سرور می‌فرستد. هکر این پیام را Intercept می‌کند و اجازه نمی‌دهد به سرور واقعی برسد:
    $$ A \rightarrow C(S): A, B$$
    
 هکر خودش را جای سرور S جا می‌زند و پیام‌های قدیمی و لو رفته‌ی نشست قبلی را برای A بازتاب می‌دهد. از آنجا که پیام با کلید $K_A$ رمز نشده، هکر می‌تواند یک کلید تصادفی دلخواه خودش ($R_{attacker}$) را با کلید $K_A$ که احتمالا از قبل با یک حمله دیگر یا نشست قدیمی به دست آورده رمز کند، یا اگر فرض کنیم هکر کلیدهای قدیمی را دارد، پیام زیر را می‌فرستد:
    $$ C(S) \rightarrow A: E(K_A, R_{attacker}), E(K_B, R_{old})$$
    
 کاربر A کلید $R_{attacker}$ را استخراج کرده، پیام جدید $M$ را با آن رمز می‌کند و به همراه تیکت قدیمی B ارسال می‌کند:
    $$ A \rightarrow C(B): E(K_B, R_{old}), E(R_{attacker}, M)$$

هکر پیام $E(R_{attacker}, M)$ را Intercept می‌کند. چون خودش کلید $R_{attacker}$ را تعیین کرده بود، به راحتی آن را رمزگشایی کرده و به متن آشکار پیغام $M$ دسترسی پیدا می‌کند. این نشان می‌دهد که پروتکل بدون استفاده از نانس و احراز اصالت دقیق سرور کاملا ناامن است.